\setAuthor{Carel Kuusk}
\setRound{piirkonnavoor}
\setYear{2019}
\setNumber{G 1}
\setDifficulty{2}
\setTopic{Kinemaatika}

\prob{Rong}
Juku tahtis teada, mitu vagunit on rongis. Selleks seisis ta perroonil esimese vaguni esiotsaga kohakuti ning ootas, kuni rong ühtlase kiirendusega sõitmist alustas.
Stopperiga mõõtes sai ta teada, et esimene vagun läks temast mööda täpselt aja $t_1=\SI{10}{s}$ jooksul ning viimane vagun aja  $t_2 = \SI{1.83}{s}$ jooksul. Mitu vagunit oli rongis?
Vagunid on identsed.\hint
Tähistades rongide arvu $n$-ga, saame kirja panna ühe vaguni möödumise aja, esimese $n-1$ vaguni möödumise aja ning kõikide vagunite aja ning tekkinud võrrandisüsteemi lahendada.\solu
Olgu rongi kiirendus $a$ ning vagunite arv $n$. Kuna vagunid on identsed, siis ühe vaguni pikkus on $s=at_1^2/2$. Kõik vagunid mööduvad Jukust aja 
aja $t_n=\sqrt{2sn/a}$ jooksul ning kõik peale viimase vaguni aja $t_{n-1}=\sqrt{2s(n-1)/a}$ jooksul. Seega 
\begin{align*}
  t_2=t_n-t_{n-1}&=\sqrt{\frac{2sn}{a}}-\sqrt{\frac{2s(n-1)}{a}}=t_1(\sqrt{n}-\sqrt{n-1})\implies\\
  \sqrt{n - 1} &= \sqrt{n} - \frac{t_2}{t_1} \implies n - 1 = n - 2\sqrt{n}\frac{t_2}{t_1} + \left(\frac{t_2}{t_1}\right)^2 \implies\\
  n &= \frac{1}{4}\left(\frac{t_2}{t_1} + \frac{t_1}{t_2}\right)^2 = 7.97 \approx 8.
\end{align*}
Teisisõnu, rongis oli 8 vagunit.\probend